\scanpage{11}
\uline{Dualité locale}. \quad $A$ anneau local noeth. \uline{complet} \uline{régulier} de dim $n$

$X=\operatorname{Spec} A$, $x$ son p\textsuperscript{t} fermé

$F$ Module variable sur $X$ (\uuline{pas néc. quasi-cohérent!})
\[
 H_x^i(F)\times\mathbb E\mathrm{xt}^{n-i}(X;F,\underline O_X)\longrightarrow H_x^n(\underline O_X)=I
 \qquad(I\ \text{module dualisant})
\]
d'où
\[
 \boxed{\operatorname{Ext}^{n-i}(X;F,\underline O_X)\longrightarrow\operatorname{Hom}(H_x^i(F),I)}
\]
Je dis que c'est un \uline{isomorphisme} \uline{pour tt} $i$, \uline{à peu de choses près}.

On est ramené à le prouver pour un faisceau $F$ de la forme $\underline O_{U,X}=i_!(\underline O_U)$, où $U$ est un ouvert de $X$ (qu'on peut supposer affine, de la forme $X_f$) et $i:U\longrightarrow X$ l'immersion canonique. La formule devient alors
\[
 \operatorname{Ext}^{n-i}(U;\underline O_{X},\underline O_X)\xrightarrow{\sim}\operatorname{Hom}(H_x^i(\underline O_{U,X}),I)
\]
i.e.
\[
 \left\{\begin{aligned}
 H_x^i(\underline O_{U,X})&=0\quad\text{si }i\ne n\\
 D(H_x^n(\underline O_{U,X}))&\approx\deleted{\operatorname{Hom}}\Gamma(U,\underline O_X)=A_f
 \end{aligned}\right.
\]
\deleted{\unsure{Ceci}} Or \deleted{\unsure{on}} si $U=X$, c'est connu, Reste $U\ne X$, on a

$\deleted{\unsure{H^0_x(\underline O}}$

1) $n=0$ on obtient une Tautologie.

2) $n=1$, il suffit de prendre $U=$ ouvert générique, et la formule devient, puisque $H^0_x(\underline O_{U,X})=0$, $H^1_x(\underline O_{U,X})=K$,
\[
 \operatorname{Hom}(K,K/A)\simeq K
\]
(autodualité bien connue!). Noter que \deleted{\unsure{pour $F$|pour}} la catégorie des $\underline O_X$-Modules est équivalente à la catégorie des triplets $(M_0,M_1,u)$, où $M_0$ est un $A$-module, $M_1$ un $K$-module [resp. $F_x$ et $F_{x'}$, où $x'$ est le pt gén.) \deleted{et} $u:M_0\longrightarrow M_1$ un $(A\longrightarrow K)$-hom. On aura
\[
 H_x^0(F)=\operatorname{Ker}u,\quad H_x^1(F)=\operatorname{Coker}u
\]
\margin{$[H^0(X,F)\simeq M_0,\ H^i(X,F)=0\ i\ne0]$\\ $\deleted{\unsure{H^i(X,\underline O_X)=\operatorname{Ext}(F,\underline O_X)}}$}
% archive sheet number: 10

